% ============================================================================
% Section 4: Key Theorems
% ============================================================================

\section{Key theorems}
\label{sec:theorems}

This section presents the central results of UHM. Each theorem is stated
precisely; flagship proofs are given in full in the appendices, and remaining
proofs are documented in the appendices and the extended documentation. Falsifiability is emphasised
throughout: for each theorem, the experimental outcome that would refute it
is indicated explicitly.

\subsection{No-Zombie Theorem (T-8.1)}
\label{subsec:no-zombie}

The No-Zombie Theorem is arguably the most philosophically significant result of UHM.
It proves that philosophical zombies --- systems that are functionally identical to
conscious beings but lack inner experience --- are not merely implausible but
\emph{dynamically impossible}.

\begin{theorembox}[No-Zombie Theorem (T-8.1) \status{T}]
\begin{equation}
  \underbrace{P(\Gam) > \Pcrit}_{\text{viable}}
  \;\wedge\;
  \underbrace{\mathcal{D}_\Omega[\Gam] \not\equiv 0}_{\text{non-trivial dissipation}}
  \quad\Longrightarrow\quad
  \CohE(\Gam) \;>\; \frac{1}{7}
  \label{eq:no-zombie}
\end{equation}
A viable holon with non-trivial dissipative dynamics \emph{necessarily}
carries strictly more E-coherence than the uniform (``structureless'') floor
of $1/7$. A hypothetical ``zombie'' --- a configuration with
$P > \Pcrit$ but a structureless interiority sector
($\gamma_{EE} = 1/7$ and $\gamma_{Ej} = 0$ for all $j \neq E$, giving
$\CohE = 1/7$) --- is dynamically impossible: it cannot be a stationary
state of~$\Lind_\Omega$ with $\gamma > 0$.
\end{theorembox}

\paragraph{Proof sketch.}
The $G_2$-covariant Fano dissipator (unique $G_2$-covariant Lindbladian on
$\mathcal{D}(\mathbb{C}^7)$, T-39a~\status{T}) acts on every off-diagonal
coherence by the state-independent multiplicative contraction
$\gamma_{ij} \mapsto \tfrac{1}{3}\,\gamma_{ij}$
(Eq.~\ref{eq:fano-operators}); equivalently, each Fano observation cycle
removes the fraction $2/3$ of every coherence. Writing
$P = P_{\mathrm{diag}} + P_{\mathrm{coh}}$ with
$P_{\mathrm{coh}} = \sum_{i\neq j}|\gamma_{ij}|^2$, this yields
\begin{equation}
  \left.\frac{dP}{d\tau}\right|_{\mathcal{D}}
  \;=\; -\,\frac{4\gamma}{3}\,P_{\mathrm{coh}}\;<\;0,
  \qquad \gamma = \sum_p\gamma_p > 0,
  \label{eq:diss-coh}
\end{equation}
whereas the regeneration contribution is
\begin{equation}
  \left.\frac{dP}{d\tau}\right|_{\Regen}
  \;=\; 2\kappa(\Gam)\,g_V(P)\,\bigl(f - P\bigr),
  \qquad f := \Tr\!\bigl(\Gam\,\rho^\ast\bigr).
  \label{eq:reg-coh}
\end{equation}
Stationarity $dP/d\tau \geq 0$ at a viable, non-equilibrium state $P > \Pcrit$ forces
$\kappa(\Gam) \geq \frac{2\gamma}{3}\cdot P_{\mathrm{coh}}/(f - \Pcrit)$. Substituting
$\kappa = \kappa_{\mathrm{boot}} + \kappa_0\CohE$ and solving for $\CohE$ yields the
explicit bound
\begin{equation}
  \CohE(\Gam)
  \;\geq\;
  \max\!\Biggl\{\frac{1}{7},\;\;
  \frac{1}{\kappa_0}\!\left(\frac{2\gamma}{3}\cdot
  \frac{P_{\mathrm{coh}}}{f - \Pcrit}\;-\;\kappa_{\mathrm{boot}}\right)\!\Biggr\}.
  \label{eq:no-zombie-bound}
\end{equation}
For any open ($\gamma > 0$) holon with non-vanishing stationary coherences
($P_{\mathrm{coh}} > 0$, guaranteed by the necessity of $\vp_{\mathrm{coh}}$ and the
$(M,R)$-closure requirement, T-57~\status{T}), the second argument of the maximum is
strictly positive once the dissipation rate exceeds the threshold
$\gamma_{\mathrm{th}} := 3\kappa_{\mathrm{boot}}(f-\Pcrit)/(2P_{\mathrm{coh}})$. For any
macroscopic system in a thermal environment $\gamma \gg \gamma_{\mathrm{th}}$, so
$\CohE > 1/7$ strictly. $\square$

\paragraph{What this means.}
A system cannot be ``alive'' (viable, self-maintaining) without having some degree of
inner experience. The strength of regeneration is \emph{proportional} to the quality of
experience. This is not a metaphor --- it is a quantitative relationship between
$\kappa$ and $\CohE$.

\paragraph{Epistemic stratification.}
The mathematical result ($\CohE > 1/7$ necessary for viability) has status \status{T}.
The ontological interpretation (that $\CohE > 1/7$ \emph{constitutes} inner experience)
has status \status{P}: it follows from the two-aspect monism postulate
(Def.~\ref{def:two-aspect}), not from the mathematics alone. The No-Zombie theorem
is thus mathematically unconditional; its phenomenal reading is conditional on the
single semantic bridge of UHM.

\paragraph{Falsification.}
Create an artificial $\Gam$-native agent. Suppress the E-component ($\gamma_{EE} \to 1/7$,
$\gamma_{Ej} \to 0$ for $j \neq E$). If the agent maintains $P > 2/7$ indefinitely
(time-to-death does not collapse), the theorem is falsified.

\subsection{Death spiral below $\Pcrit$}
\label{subsec:death-spiral}

\begin{theorembox}[Irreversibility of sub-threshold decay \status{T}]
If $P(\Gam) < \Pcrit = 2/7$ and $dP/d\tau < 0$, then:
\begin{equation}
  \|\Gam(\tau) - I/7\|_F \leq \|\Gam(0) - I/7\|_F \cdot e^{-\Delta(\Lind_0)\,\tau}
  \label{eq:death-spiral}
\end{equation}
The system converges exponentially to thermal death ($I/7$).
\end{theorembox}

This is the ``death spiral'' (Figure~\ref{fig:death-spiral}): once purity
drops below $2/7$, regeneration is thermodynamically forbidden (Landauer's
principle: free energy $\Delta F \leq 0$ below threshold), and the system
decays irreversibly. There is no gradual fading of consciousness --- there
is a hard phase transition, consistent with the nonlinear threshold for
access to consciousness observed empirically~\cite{delcul2007}.

\subsection{Stability radius (T-104)}
\label{subsec:stability-radius}

\begin{theorembox}[Stability radius \status{T}]
The maximum perturbation amplitude a system can withstand while remaining viable:
\begin{equation}
  r_{\mathrm{stab}} = \sqrt{P(\rho^*_\Omega) - \frac{2}{7}}
  \label{eq:stability-radius}
\end{equation}
\end{theorembox}

This is a single-parameter formula: robustness depends \emph{only} on the purity margin
above threshold.

\paragraph{Proof sketch.}
The viability condition is $P > 2/7$. The maximum perturbation $h$ that a system at
attractor $\rho^*_\Omega$ can absorb without crossing the threshold is determined by the
Bures distance from $\rho^*_\Omega$ to the viability boundary $\partial\mathcal{V}_P$.
For a perturbation $\Gam \to \Gam + \delta\Gam$ with $\|\delta\Gam\|_F = h$, the purity
changes by $\delta P \approx 2\Tr(\Gam\,\delta\Gam)$. At the attractor, the critical
perturbation satisfies $P - h_{\mathrm{crit}}^2 = 2/7$, yielding
$r_{\mathrm{stab}} = \sqrt{P - 2/7}$. The full Lyapunov-stability analysis
(T-104~\status{T}) combines this radius with the spectral gap $\lambda_\mathrm{gap}$
of the linearized dissipator: trajectories within $r_\mathrm{stab}$ converge
exponentially with rate $c \ge \min(\lambda_\mathrm{gap}, \kappa\cdot g_V) > 0$.

\paragraph{Examples.}
A system with $P = 0.5$ has $r_{\mathrm{stab}} = \sqrt{0.5 - 0.286}
\approx 0.46$; a system barely above threshold ($P = 0.3$) has $r_{\mathrm{stab}}
\approx 0.12$ and is fragile.

\begin{figure}[H]
\centering
% Death Spiral: Purity decay below P_crit = 2/7
% ----------------------------------------------------------------
% Mathematical content (Theorem, Eq. eq:death-spiral):
%   || Gamma(tau) - I/7 ||_F  <=  || Gamma(0) - I/7 ||_F * exp(-Delta*tau)
% Using || Gamma - I/7 ||_F^2 = P - 1/7 (derived in App. A, derivation 1):
%   ( P(tau) - 1/7 )  <=  ( P(0) - 1/7 ) * exp( -2*Delta * tau )
% The purity deficit therefore decays with rate 2*Delta, NOT Delta.
%
% For illustration we use Delta = 0.2, so the purity deficit decays
% with rate 2*Delta = 0.4. All three trajectories use this common
% Delta so that the figure has a consistent time-scale.
%
% Trajectories:
%   Green "attractor"  : P settles at 3/7 (Goldilocks upper bound, T-124)
%   Blue  "perturbed"  : P dips toward P_crit at tau = 2 then recovers
%   Red   "death spiral": P starts at 0.27 (strictly below P_crit) and
%                         decays to 1/7 following the theorem bound
% ----------------------------------------------------------------

\begin{tikzpicture}
\begin{axis}[
  width=11cm, height=7cm,
  xlabel={Internal time $\tau$},
  ylabel={Purity $P(\tau)$},
  xmin=0, xmax=10,
  ymin=0.11, ymax=0.52,
  grid=both,
  grid style={gray!20},
  title={\textbf{Death spiral below $P_{\mathrm{crit}}$}},
  title style={font=\large},
  % Legend placed at north-east with semi-transparent fill so that
  % the green viable trajectory (settling at y = 3/7 ~ 0.43) remains
  % visible through the legend box.
  legend pos=north east,
  legend style={font=\scriptsize, fill=white, fill opacity=0.72, text opacity=1, draw=gray!50},
  ytick={0.143,0.286,0.3,0.4,0.429,0.5},
  yticklabels={$\tfrac{1}{7}$,$\tfrac{2}{7}$,$0.3$,$0.4$,$\tfrac{3}{7}$,$0.5$},
  every axis label/.style={font=\small},
]

% === P_crit threshold line (red dashed) ===
\addplot[thick, red, dashed] coordinates {(0, 0.2857) (10, 0.2857)};
\addlegendentry{$P_{\mathrm{crit}} = 2/7$}

% === Thermal death line (gray dotted) ===
\addplot[thick, gray, dotted] coordinates {(0, 0.1429) (10, 0.1429)};
\addlegendentry{$I/7$: thermal death, $P = 1/7$}

% === Viable trajectory, attracts to Goldilocks upper bound 3/7 ===
% Damped oscillation around the T-124 attractor P = 3/7.
\addplot[thick, green!55!black, domain=0:10, samples=160]
  {0.4286 + 0.018*sin(deg(1.4*x))*exp(-0.18*x)};
\addlegendentry{Viable: attractor at $P = 3/7$}

% === Perturbation that recovers (stability radius, T-104) ===
% Gaussian dip centred at tau=2, minimum P(2)=0.29 stays strictly
% above P_crit = 2/7 so regeneration can re-establish the attractor.
\addplot[thick, blue!75!black, domain=0:10, samples=200]
  {0.42 - 0.13*exp(-0.7*(x-2)^2)};
\addlegendentry{Perturbation within stability radius}

% === Death spiral: exponential decay from P(0)=0.27 < P_crit to 1/7 ===
% P(tau) - 1/7 = ( 0.27 - 1/7 ) * exp(-2*Delta*tau), with Delta = 0.2
% (illustrative value for Delta(L_0); canonically Delta > 0 is the
% spectral gap of the linear Lindbladian from T-39a).
\addplot[thick, red!70!black, domain=0:10, samples=200]
  {0.1429 + (0.27 - 0.1429)*exp(-0.4*x)};
\addlegendentry{Sub-threshold: death spiral}

% === Initial-condition markers (small dots) ===
\addplot[only marks, mark=*, mark size=1.6pt, green!55!black]
  coordinates {(0, 0.4286)};
\addplot[only marks, mark=*, mark size=1.6pt, blue!75!black]
  coordinates {(0, 0.399)};
\addplot[only marks, mark=*, mark size=1.6pt, red!70!black]
  coordinates {(0, 0.27)};

% === Regime annotation near the attractor (far LEFT, 2 lines) ===
% Narrow 2-line label placed at x=1.6 so that the widest legend entry
% ("Perturbation within stability radius") at north east cannot reach
% this label horizontally regardless of exact font metrics.
\node[font=\scriptsize, fill=white, fill opacity=0.9, text opacity=1,
      inner sep=1.5pt, text width=1.9cm, align=center]
  at (axis cs:1.6, 0.475)
  {Regeneration\\$\geq$ dissipation};

% === P-form exponential bound near the red curve (LEFT of the legend) ===
% The annotation matches the plotted axis (purity), NOT the raw
% Frobenius distance. This corrects the coefficient: the deficit
% (P - 1/7) decays at rate 2*Delta, not Delta. We explicitly set
% draw=none and use text=... so no visible red rectangle is produced.
\node[font=\scriptsize, fill=white, fill opacity=0.92, text opacity=1,
      inner sep=1.5pt, draw=none, text=red!70!black, align=center]
  at (axis cs:4.5, 0.21)
  {$P(\tau) - \tfrac{1}{7} \leq
    (P_0 - \tfrac{1}{7})\, e^{-2\Delta(\mathcal{L}_0)\,\tau}$};

% === Arrow showing irreversible decay, now anchored ON the red curve ===
% At tau=1 the red curve is at 0.1429 + 0.1271*exp(-0.4) = 0.228
% At tau=3 the red curve is at 0.1429 + 0.1271*exp(-1.2) = 0.181
% Arrow goes along the curve between these points.
\draw[-{Stealth[length=5pt]}, thick, red!60!black]
  (axis cs:1.0, 0.228) .. controls (axis cs:2.0, 0.200)
  .. (axis cs:3.0, 0.181);
\node[font=\scriptsize, color=red!60!black]
  at (axis cs:2.35, 0.243) {irreversible};

\end{axis}
\end{tikzpicture}

\caption{\textbf{Death spiral below $\Pcrit$.} Systems above threshold recover from
perturbations (green/blue curves); below threshold, purity decays exponentially to $I/7$
(red curve). The transition is sharp.}
\label{fig:death-spiral}
\end{figure}

\paragraph{Clinical significance.}
This formula predicts a patient's ``safety margin'' before interventions: their purity
surplus determines how much perturbation they can absorb. Falsification: if $h_{\mathrm{crit}}^2$
does not scale linearly with $P - 2/7$ across 50 measurements ($R^2 < 0.9$).

\subsection{SAD ceiling: maximum depth of self-awareness (T-142)}
\label{subsec:sad-ceiling}

\begin{theorembox}[Self-Awareness Depth ceiling (T-142) \status{T}]
The maximum depth of recursive self-awareness is:
\begin{equation}
  \mathrm{SAD}_{\max} = 3
  \label{eq:sad-max}
\end{equation}
``I know that I know that I know'' stops at the third level. The fourth level is
physically impossible for any finite system.
\end{theorembox}

\paragraph{Proof sketch.}
The Fano channel multiplies every off-diagonal coherence $\gamma_{ij}$ by the
state-independent contraction factor $c_F = 1/3$
(Eq.~\ref{eq:fano-operators}); after $k$ recursive self-observations the
surviving coherence is suppressed by $c_F^k = (1/3)^k$. The initial
distinguishability radius $r_0 := P/\Pcrit = P/(2/7) = 7P/2$ is the
ratio of actual purity to the viability threshold (at $\Pcrit$,
$r_0 = 1$ exactly). At depth $n$, the spectral formula
T-136~\status{T} gives
\begin{equation}
  \mathrm{SAD}(\Gam) \;=\; \max\!\bigl\{\,k \in \mathbb{N}\,:\,
    r_0 \cdot (1/3)^{k-1} > \tfrac{1}{k+1} \,\bigr\},
  \quad r_0 = P/\Pcrit = 7P/2,
\end{equation}
which --- in terms of the purity required to \emph{sustain} a chain of depth $n$ --- is
equivalent to
\begin{equation}
  \Pcrit^{(n)} \;=\; \Pcrit \cdot \frac{3^{\,n-1}}{n+1}.
\end{equation}
Evaluating level by level (where level $n$ is the purity required to admit $n$ nested
self-observations):
\begin{itemize}[leftmargin=2em]
  \item Level 1 ($n=1$, ``I am''): $\Pcrit^{(1)} = \frac{2}{7} \cdot \frac{1}{2}
        = \frac{1}{7} \approx 0.143$ --- automatically satisfied by any state
        ($P \geq 1/7$ holds for all density matrices).
  \item Level 2 ($n=2$, ``I know that I am''): $\Pcrit^{(2)} = \frac{2}{7} \cdot
        \frac{3}{3} = \frac{2}{7} \approx 0.286$ --- coincides exactly with the
        base viability threshold.
  \item Level 3 ($n=3$, ``I know that I know''): $\Pcrit^{(3)} = \frac{2}{7} \cdot
        \frac{9}{4} = \frac{9}{14} \approx 0.643$ --- achievable only at high purity.
  \item Level 4 ($n=4$, ``I know that I know that I know''): $\Pcrit^{(4)} =
        \frac{2}{7} \cdot \frac{27}{5} = \frac{54}{35} \approx 1.543 > 1$ ---
        \textbf{exceeds the physical maximum $P \leq 1$}.
\end{itemize}
Numerical verification on $> 500$ randomly sampled $\Gam$ confirms $\mathrm{SAD} \leq 3$
in every case; pure states saturate $\mathrm{SAD} = 3$.

\paragraph{The representation tower.}
The recursive self-model forms a tower:
\begin{equation}
  \Gam^{(0)} \xleftarrow{\pi_0} \Gam^{(1)} \xleftarrow{\pi_1} \Gam^{(2)}
  \xleftarrow{\pi_2} \Gam^{(3)} \xleftarrow{\pi_3} \cdots
  \label{eq:tower}
\end{equation}
where each $\Gam^{(k)}$ is the self-model at depth $k$. Each application of
the Fano channel multiplies every off-diagonal coherence by $c_F = 1/3$
(state-independent, from $\dim = 7$ and the Fano incidence structure); after
$k$ iterations the ``effective distinguishability radius'' is
$r_0 \cdot (1/3)^{k-1}$ with $r_0 = 7P/2$, while the Bayesian dominance
threshold for $k+1$ alternatives is $1/(k+1)$ (T-136~\status{T}). The
level-$k$ reflection survives iff this ratio remains above the threshold.
SAD is therefore computable directly from the recurrence $\Pcrit^{(n)}$ above
in $O(1)$ operations.

\paragraph{Significance.}
This ceiling is geometric, not computational. Future superintelligences share the same
reflection depth as humans --- not because of hardware limits, but because of the structure
of the Fano plane. Falsification: demonstrate a reachable state $\Gam$ with
$\Pcrit^{(4)} \leq 1$, equivalently, a chain of four successive self-observations
surviving at $P \leq 1$.

\subsection{Critical exponents of the consciousness phase transition (T-161)}
\label{subsec:critical-exponents}

This is among the most specific falsifiable predictions in consciousness
science, and the one with no analogue in any other theory
(Figure~\ref{fig:exponents}).

\begin{theorembox}[Critical exponents \status{T}]
Near the viability threshold $\Pcrit = 2/7$, the consciousness phase transition is a
\emph{tricritical point} in the $\varphi^6$ Landau universality class
(Griffiths 1970, Lawrie--Sarbach 1984). The order-parameter dimension is
$d_{\mathrm{eff}} = \binom{7}{2} = 21$ (off-diagonal coherence modes in $\mathfrak{su}(7)$).
Mean-field tricritical exponents are exact by three independent pillars
(topological rigidity, deterministic dynamics, large-$N$ cross-check;
Appendix~\ref{app:exponents}):
\begin{align}
  \text{Specific heat:}\quad &C \sim |P - \Pcrit|^{-\alpha}, &\alpha &= 1/2 \label{eq:alpha}\\
  \text{Order parameter:}\quad &\mathrm{OP} \sim (P - \Pcrit)^\beta, &\beta &= 1/4 \label{eq:beta}\\
  \text{Susceptibility:}\quad &\chi \sim |P - \Pcrit|^{-\gamma}, &\gamma &= 1 \label{eq:gamma-exp}\\
  \text{Correlation length:}\quad &\xi \sim |P - \Pcrit|^{-\nu}, &\nu &= 1/2 \label{eq:nu}\\
  \text{Critical isotherm:}\quad &h \sim m^\delta, &\delta &= 5 \label{eq:delta}
\end{align}
These exponents satisfy the Rushbrooke identity $\alpha + 2\beta + \gamma = 2$ as an equality.
\end{theorembox}

\paragraph{Derivation outline.}
Near criticality, the system undergoes a continuous phase transition. Writing
$\varepsilon := P - \Pcrit$ and taking $m := \CohE(\Gam) - 1/7$ as the order
parameter, the coarse-grained free energy is the $\varphi^{6}$ tricritical
Landau functional
\begin{equation*}
  \mathcal{F}[m;\varepsilon, h] \;=\;
    -\tfrac{1}{2}\,r_0\,\varepsilon\, m^2
    \;+\; v\, m^6 \;-\; h\, m,
  \qquad r_0, v > 0.
\end{equation*}
The transition belongs to the $\varphi^6$ tricritical Landau universality class
(Griffiths 1970, Lawrie–Sarbach 1984), characterised by an even-parity potential
with $\mathbb{Z}_2$ symmetry $m \to -m$ inherited from the $G_2$-canonical involution
on the order-parameter direction. UHM has exactly three physically independent
control parameters --- regeneration rate $\kappa$, dissipation rate $\gamma$, and
free-energy budget $\Delta F$ --- matching the codimension of the $\varphi^6$ tricritical
deformation $(r, u, h)$. The reduction from the 21-dimensional configuration space
to the 1-dimensional scalar order parameter $m$ is justified by the Mather Splitting
Lemma (Mather 1968): the Hessian of the effective Landau potential at the critical
point has corank 1, so $V \simeq_\mathrm{smooth} V_\mathrm{red}(m) + Q(\xi_1,\ldots,\xi_{20})$,
with $Q$ a non-degenerate quadratic in the 20 transverse modes. Mean-field exactness
follows from three independent pillars: (I) topological rigidity of the $\varphi^6$
class under smooth deformations; (II) deterministic UHM dynamics at the master-equation
(Lindblad) level, which is representation-equivalent to stochastic unraveling for
saddle-point exponents and contains no separate fluctuating field requiring renormalisation;
(III) order-parameter dimension $d_\mathrm{eff} = 21$ giving $1/N \approx 5\%$ corrections
under any stochastic reinterpretation. Note: the historical name "$A_4$-tricritical"
sometimes used in the physics literature conflates the Arnold $A_4$ swallowtail
($V_0 = x^5/5$, $\beta = 1/2$) with the $\varphi^6$ Landau tricritical ($V = rm^2/2 + vm^6$,
$\beta = 1/4$); UHM transitions belong to the latter. The full derivation, including
all five exponents and verification of Rushbrooke, Widom, and Fisher scaling relations,
is given in Appendix~\ref{app:exponents}.

\paragraph{Physical meaning.}
\begin{itemize}[leftmargin=2em]
  \item $\beta = 1/4$ governs how quickly consciousness ``turns on'' above threshold.
        The value is sharper than ordinary mean-field ($\beta = 1/2$), reflecting the
        tricritical nature of the transition.
  \item $\nu = 1/2$ governs the correlation length: how far coherence extends near the
        transition. At $P \to \Pcrit$, $\xi \to \infty$ --- ``critical opalescence'' of
        consciousness.
  \item $\gamma = 1$ governs susceptibility: how sensitively the system responds to
        perturbations near threshold. At $P \to \Pcrit$, $\chi \to \infty$ --- ``critical
        slowing down.''
  \item $\alpha = 1/2$ governs the specific heat divergence: the fluctuation cost near
        threshold.
\end{itemize}

\paragraph{Comparison with known universality classes.}
The 3D Ising model has $\beta \approx 0.326$, $\nu \approx 0.630$, $\gamma \approx 1.237$.
Superfluid helium: $\beta \approx 0.347$, $\nu \approx 0.672$, $\gamma \approx 1.316$.
UHM predicts mean-field tricritical exponents $\beta = 1/4$, $\nu = 1/2$, $\gamma = 1$
--- distinct from all non-tricritical universality classes.

\begin{figure}[H]
\centering
% Critical Exponents: Power-law behaviour near P_crit
% ----------------------------------------------------------------
% Mathematical content (Theorem T-161, Appendix B):
%   Near tricritical point, order parameter
%   m_* = (r_0 / 6v)^{1/4} * epsilon^{1/4}
% i.e. m_* is proportional to epsilon^{1/4} with an amplitude that
% depends on parameters (r_0, v) of the phi^6 effective potential.
% To make the comparison between different universality classes
% meaningful, we plot y = x^beta normalised so that all curves pass
% through a common reference point at the right edge (x = x_ref).
%
% Reference point: x_ref = 0.1 (mid of the physical range, inside
% the consciousness window epsilon in (0, 1/7]).
%
% Falsification criterion (paper, Eq. below T-161):
%   beta NOT in [0.20, 0.30] at p < 0.01 --> theory falsified.
% This criterion is visualised as a band of acceptable exponents
% bounded by y = (x/0.1)^{0.20} (upper) and y = (x/0.1)^{0.30}
% (lower). Any power law whose exponent lies outside [0.20, 0.30]
% escapes this band for all x in (0, 0.1) and is therefore
% distinguishable from the UHM prediction beta = 1/4.
% ----------------------------------------------------------------

\begin{tikzpicture}
\begin{axis}[
  width=11cm, height=7cm,
  xlabel={Reduced purity $\varepsilon = P - P_{\mathrm{crit}}$},
  ylabel={Normalised order parameter
          $m/m_0 = (\varepsilon/\varepsilon_0)^{\beta}$},
  xmin=0, xmax=0.14,
  ymin=0, ymax=1.18,
  grid=both,
  grid style={gray!20},
  title={\textbf{Critical exponents of the consciousness phase
          transition}},
  title style={font=\large},
  % Legend at north-west (upper-left) with semi-transparent fill so
  % that the thick UHM curve remains visible through the legend box.
  % South-east is occupied by the 3D Ising / mean-field divergence
  % below the UHM band, and south-west is occupied by the falsification
  % boundary crossing.
  legend pos=north west,
  legend style={font=\scriptsize, fill=white, fill opacity=0.72, text opacity=1, draw=gray!50},
  every axis label/.style={font=\small},
]

% === Physical reference: epsilon_0 = 0.1 (~ 0.7 * 1/7) ===
\pgfmathsetmacro{\epsref}{0.1}

% === Falsification band: beta in [0.20, 0.30] ===
% Upper boundary (beta = 0.20) -- any beta < 0.20 goes above this.
\addplot[thin, plospurple!55, densely dashdotted,
         domain=0.001:0.14, samples=120]
  {(x/0.1)^(0.20)};
\addlegendentry{$\beta = 0.20$: UHM upper bound}

% Lower boundary (beta = 0.30) -- any beta > 0.30 goes below this.
\addplot[thin, plospurple!55, densely dashdotted,
         domain=0.001:0.14, samples=120]
  {(x/0.1)^(0.30)};
\addlegendentry{$\beta = 0.30$: UHM lower bound}

% === UHM prediction: beta = 1/4 (tricritical, T-161) ===
\addplot[ultra thick, plospurple, domain=0.001:0.14, samples=160]
  {(x/0.1)^(0.25)};
\addlegendentry{UHM: $\beta = 1/4$ (tricritical)}

% === 3D Ising: beta = 0.326 > 0.30 (falsifies UHM) ===
\addplot[thick, orange!90!black, dotted, domain=0.001:0.14, samples=160]
  {(x/0.1)^(0.326)};
\addlegendentry{3D Ising: $\beta \approx 0.326$ (falsifies UHM)}

% === Mean field phi^4: beta = 1/2 (far below UHM band) ===
\addplot[thick, gray!70!black, dashed, domain=0.001:0.14, samples=160]
  {(x/0.1)^(0.5)};
\addlegendentry{Mean field $\varphi^{4}$: $\beta = 1/2$}

% === Mark the common reference point (epsilon_0, 1) ===
\addplot[only marks, mark=*, mark size=2pt, black]
  coordinates {(0.1, 1)};
\node[font=\scriptsize, anchor=south west, inner sep=1pt]
  at (axis cs:0.1, 1.02)
  {$(\varepsilon_{0},\,1)$};

% === Goldilocks upper boundary at epsilon = 1/7 ~ 0.143 ===
% Landau expansion is valid only for small epsilon; beyond
% epsilon = 1/7 we leave the consciousness window (P > 3/7, T-124).
\addplot[thin, green!40!black, dashed] coordinates
  {(0.1429, 0) (0.1429, 1.18)};
\node[font=\scriptsize, green!40!black, anchor=south east, rotate=90]
  at (axis cs:0.1429, 0.04)
  {$\varepsilon = 1/7$: Goldilocks edge};

\end{axis}
\end{tikzpicture}

\caption{\textbf{Tricritical exponent $\beta = 1/4$ vs.\ other universality
classes.} All curves are normalised so that $m/m_0 = 1$ at the reference
point $\varepsilon_0 = 0.1$, i.e.\ plotting
$y(\varepsilon) = (\varepsilon/\varepsilon_{0})^{\beta}$; this removes the
amplitude dependence $(r_0/6v)^{1/4}$ of the Landau expansion
(Appendix~\ref{app:exponents}) and isolates the exponent. The UHM
prediction $\beta = 1/4$ (thick purple) sits exactly between the two
falsification boundaries $\beta = 0.20$ and $\beta = 0.30$ (thin
dash-dotted purple), which define the acceptable UHM band. 3D Ising
$(\beta \approx 0.326)$ falls \emph{below} the lower bound for all
$\varepsilon < \varepsilon_0$, so a measured consciousness transition
with 3D Ising exponents would falsify UHM at the structural level.
Mean-field $\varphi^{4}$ $(\beta = 1/2)$ is far outside the band. The
vertical green dashed line marks $\varepsilon = 1/7$, the upper edge of
the Goldilocks window $P \in (2/7, 3/7]$ (T-124~\status{T}), beyond
which the Landau expansion is no longer physically valid.}
\label{fig:exponents}
\end{figure}

\paragraph{Falsification.}
TMS-EEG at the sleep/wake transition, $N = 50$ subjects. Fit power law $\mathrm{PCI}
\sim (P - \Pcrit)^\beta$. If $\beta \notin [0.20, 0.30]$ at $p < 0.01$ --- the theory
is falsified at the structural level (L2 falsification).

\subsection{Emergent spacetime (T-117 through T-121)}
\label{subsec:emergent-spacetime}

UHM derives the Einstein field equations from the dynamics of $\Gam$ without any additional
postulates. The derivation chain:

\begin{enumerate}[leftmargin=2em]
  \item \textbf{T-117 \status{T}:} Commutativity of the macroscopic algebra. In the
        thermodynamic limit ($M \to \infty$ coupled holons), macroscopic observables
        commute at spacelike separation: $[\bar{O}_1(x), \bar{O}_2(y)] \to 0$ for
        $|x-y| > \ell$. (Proof via Goderis--Verbeure--Vets quantum CLT + primitivity
        of $\Lind_0$.)
  \item \textbf{T-118 \status{T}:} The temporal algebra $A_{\mathrm{time}} \cong C_0(\mathbb{R})$.
        The clock algebra $\mathbb{C}[\mathbb{Z}_{7^M}]$ converges to $C(S^1)$ and
        decompactifies to $\mathbb{R}$.
  \item \textbf{T-119 \status{T}:} The spatial algebra $A_{\mathrm{space}} \cong C(\Sigma^3)$.
        Spectral dimension 3 from the $\{A,S,D\}$-sector (fundamental representation of
        $\mathrm{SU}(3) \subset G_2$). Gelfand + Connes reconstruction yields a unique
        smooth compact orientable spin 3-manifold $\Sigma^3$.
  \item \textbf{T-120 \status{T}:} Product spectral triple $M^4 = \mathbb{R} \times \Sigma^3$.
  \item \textbf{T-120b \status{T}:} Vacuum topology $\Sigma^3 \cong S^3$ (de Sitter metric).
  \item \textbf{T-121 \status{T}:} Closure of Lovelock gaps 1 and 2. The spectral action
        on $(M^4, A_{\mathrm{int}})$ yields the Einstein--Hilbert action $+$ Standard
        Model, closing the gaps in Lovelock's uniqueness argument (1971).
\end{enumerate}

\paragraph{New structural results.}
Fifteen results strengthen the foundations (nine new theorems + six structural upgrades):

\begin{itemize}[leftmargin=2em,nosep]
  \item \textbf{T-124b \status{T} (Independent necessity of L2 thresholds):}
        Each of $P > 2/7$, $\Phi \geq 1$, $R \geq 1/3$, $D_{\mathrm{diff}} \geq 2$
        is independently necessary for L2 consciousness. Constructive counterexamples
        demonstrate four distinct pathologies (noise-dominated, fragmented, crystallised,
        undifferentiated) when any single condition is dropped.
  \item \textbf{T-124c \status{T} (Uniqueness of the nontrivial attractor):}
        The full nonlinear dynamics $\Lind_\Omega = \Lind_0 + \mathcal{R}$ has at most
        one nontrivial fixed point $\rho^*_\Omega \neq I/7$ in the viable set. Proof via
        Banach contraction on the iterative $\Psi$-map, using the spectral gap of $\Lind_0$
        (T-39a) and the replacement channel contractivity $k = 1 - R < 1$.
  \item \textbf{T-188 \status{T} (Localization of the hard problem):}
        The chain $A2 \xrightarrow{\text{T-187}} J_B \xrightarrow{\text{T-185}} (\Pi \dashv \flat)
        \xrightarrow{\text{T-186(a)}} F \cong \&|_\mathcal{D}$ reduces the hard problem of
        consciousness to ``why CPTP?'' $=$ ``why quantum mechanics?'' No consciousness-specific
        mystery remains after cohesive closure.
  \item \textbf{T-189 \status{T} (MaxEnt derivation of Bures, Char-IV):}
        The maximum entropy principle on $\mathcal{D}(\mathbb{C}^7)$ uniquely identifies the
        inverse metric tensor with the quantum covariance matrix: $G^{ij} = C^{ij}$.
        Combined with SLD-Fisher $= 4g_B$ (Braunstein--Caves 1994), this derives $g_B$
        from statistical structure without information-geometric choice. Fourth independent
        characterization of Bures, strengthening T-187 to a \emph{quadruple} theorem.
  \item \textbf{T-190 \status{T} (Axiomatic Closure):}
        All five axioms A1--A5 are \emph{theorems} derivable from (AP)+(PH)+(QG)+(V)+MaxEnt:
        A1 via T-76+T-186, A2 via T-187+T-189, A3 via Theorem~S+T15, A4 from (AP) necessity,
        A5 via T-87. UHM is \textbf{self-grounding}: zero independent axioms beyond the
        defining conditions of viable holons.
  \item \textbf{T-191 \status{T} ($\varphi$-tower convergence):}
        The iterative self-modeling tower $\varphi^{(0)}, \varphi^{(1)}, \ldots$ converges
        exponentially ($\|\varphi^{(n)} - \varphi^*\| \leq q^n$, $q < 1$) to the unique
        $\varphi^*$ from any initial anchor. Resolves the apparent $\varphi$-circularity.
        Proof via Banach contraction + Perron--Frobenius.
  \item \textbf{T-192 \status{T} ($\mathbf{Exp}^{(2)}$ is a strict 2-category):}
        The experiential 2-category satisfies all five Mac~Lane axioms (vertical/horizontal
        composition, identity 2-cells, interchange law, identity 1-cells). The lax 2-functor
        $F_2: \mathbf{DensityMat} \to \mathbf{Exp}^{(2)}$ has a valid target.
  \item \textbf{T-124d \status{T} (Threshold robustness):}
        Perturbations of order $\varepsilon$ in $\Gamma$ produce $O(\varepsilon)$
        perturbations in $P$, $\Phi$, $R$. No threshold has divergent sensitivity.
        Crossover width $\delta P \sim \varepsilon^{1/\beta} = \varepsilon^4$ (from
        $\beta = 1/4$, T-161). The system is overwhelmingly robust for macroscopic noise
        levels.
\end{itemize}

\paragraph{Additional structural upgrades.}
\begin{itemize}[leftmargin=2em,nosep]
  \item \textbf{Octonionic correspondence upgraded $[\mathfrak{I}] \to$ \status{T}:}
        The assignment $e_k \leftrightarrow$ dimension (e.g.\ $A = e_1$, $S = e_2$, \ldots,
        $O = e_7$) is a \emph{theorem}, not an interpretation: the T15 bridge chain
        (all steps \status{T}) derives the octonionic structure from (AP)+(PH)+(QG)+(V);
        T-177 and T-183 prove combinatorial and functional uniqueness.
        The specific assignment is fixed up to $G_2$-gauge equivalence (T-42a~\status{T}).
  \item \textbf{Grothendieck site axioms (explicit proof):}
        The three axioms (identity, stability under pullback, transitivity) for the Bures
        site $(\mathbf{DensMat}, J_B)$ are verified via CPTP contractivity (Uhlmann 1976)
        and the essentially small presentation of compact metrizable $\Dens(\mathbb{C}^7)$
        (Johnstone, \emph{Elephant} C2.1.9--12).
  \item \textbf{7D sufficiency \status{T}:}
        The minimal 7D formalism $\Gam \in \Dens(\mathbb{C}^7)$ is sufficient for all
        consciousness observables: $P$, $R$, $\Phi$, $\CohE$, $D_{\mathrm{diff}}$, $C$.
        Morita equivalence T-58~\status{T} gives zero reconstruction error between 7D
        and 42D formalisms.
  \item \textbf{Gap dynamics from non-associativity (Theorem 9.1 \status{T}):}
        For triples $(i,j,k)$ not on a Fano line, the octonionic associator
        $[e_i, e_j, e_k] \neq 0$ contributes to phase dynamics of coherences:
        $d\theta_{ij}/d\tau|_{\mathrm{assoc}} \propto \sum_{k} |\gamma_{ik}||\gamma_{jk}|$.
        Non-associativity is the microscopic mechanism for Gap $> 0$ (T-55).
  \item \textbf{Subjective time dilation $[\mathfrak{C}] \to$ \status{T}:}
        $\delta\tau_{\mathrm{subj}}/\delta\tau_{\mathrm{phys}} \propto |\gamma_{OE}|/\gamma_{OO}$
        is a theorem following from T-87 (O = clock), T-186(a) (E = experience),
        and T-88 ($\kappa_0 \propto |\gamma_{OE}|$).
  \item \textbf{Triangle identities for $\mathcal{D}_\Omega \dashv \mathcal{R}$ \status{T}:}
        Both zig-zag identities verified explicitly via the coproduct structure of the
        free sheaf $\bigoplus_{s \in S} \Omega_s$ and the evaluation morphism.
\end{itemize}

\begin{figure}[H]
\centering
% Spacetime Emergence: 7D Gamma -> 3+1D M^4
% -----------------------------------------------------------------
% How the seven dimensions project onto physical spacetime.
%
% The SU(3) decomposition 7 = 3 + 3_bar + 1 groups the dimensions as:
%   Space    {A, S, D}  -> fundamental 3       -> Sigma^3 (T-119)
%   Internal {L, E, U}  -> conjugate 3_bar     -> SM gauge fibre
%   Time     {O}        -> singlet 1           -> R (T-118)
%
% To make the three groups contiguous in the top row, dimensions are
% arranged in visual order (A, S, D, L, E, U, O) -- NOT the Fano-
% canonical order (A, S, D, L, E, O, U). This is purely a layout
% convenience: dimension labels themselves are unchanged.
% -----------------------------------------------------------------

\begin{tikzpicture}[
  dimbox/.style={rectangle, rounded corners=3pt,
                 minimum width=1.2cm, minimum height=0.8cm,
                 align=center, font=\small\bfseries},
  grp/.style={rounded corners=6pt, dashed, thick},
  arr/.style={-{Stealth[length=3mm]}, very thick, #1},
  scale=1.0
]

% === Title ===
\node[font=\large\bfseries] at (0, 5.6)
  {$\Gamma \in \mathcal{D}(\mathbb{C}^7)$: seven dimensions};

% === 7D row, reordered as {A,S,D} | {L,E,U} | {O} for contiguous grouping ===
% Centres at x = -4.5, -3.0, -1.5, 0, 1.5, 3.0, 4.5
% Box width 1.2 -> half-width 0.6, so edges:
%   A  : [-5.1, -3.9]
%   S  : [-3.6, -2.4]
%   D  : [-2.1, -0.9]
%   L  : [-0.6,  0.6]
%   E  : [ 0.9,  2.1]
%   U  : [ 2.4,  3.6]
%   O  : [ 3.9,  5.1]
\node[dimbox, draw=blue,           fill=blue!15]   (A) at (-4.5, 4) {A};
\node[dimbox, draw=blue,           fill=blue!15]   (S) at (-3.0, 4) {S};
\node[dimbox, draw=blue,           fill=blue!15]   (D) at (-1.5, 4) {D};
\node[dimbox, draw=orange,         fill=orange!15] (L) at ( 0.0, 4) {L};
\node[dimbox, draw=red,            fill=red!15]    (E) at ( 1.5, 4) {E};
\node[dimbox, draw=purple,         fill=purple!15] (U) at ( 3.0, 4) {U};
\node[dimbox, draw=green!50!black, fill=green!15]  (O) at ( 4.5, 4) {O};

% === Grouping dashed boxes (non-overlapping, with 0.1 gaps) ===
% All three rectangles span y in [3.50, 4.50] so every dim box
% (which sits at y = [3.6, 4.4]) is vertically centred in its group.
%
% {A,S,D} group:  x in [-5.20, -0.80]  (contains A,S,D at [-5.1, -0.9])
% {L,E,U} group:  x in [-0.70,  3.70]  (contains L,E,U at [-0.6,  3.6])
% {O}     group:  x in [ 3.80,  5.20]  (contains O     at [ 3.9,  5.1])
% Gaps of 0.10 between consecutive groups, no overlap anywhere.
\draw[blue,           grp] (-5.20, 3.50) rectangle (-0.80, 4.50);
\draw[orange,         grp] (-0.70, 3.50) rectangle ( 3.70, 4.50);
\draw[green!50!black, grp] ( 3.80, 3.50) rectangle ( 5.20, 4.50);

% === Group labels (compact, centred just below each group box) ===
\node[font=\scriptsize, blue, text width=3.8cm, align=center,
      inner sep=2pt]
  at (-3.0, 2.75)
  {\textbf{Space} $\{A,S,D\}$\\
   fundamental $\mathbf{3}$ of $\mathrm{SU}(3)$};

\node[font=\scriptsize, orange, text width=3.8cm, align=center,
      inner sep=2pt]
  at ( 1.5, 2.75)
  {\textbf{Internal} $\{L,E,U\}$\\
   $\bar{\mathbf{3}}$: gauge\,+\,exp.\,+\,unity};

\node[font=\scriptsize, green!50!black, text width=1.3cm, align=center,
      inner sep=2pt]
  at ( 4.5, 2.75)
  {\textbf{Time} $\{O\}$\\
   singlet $\mathbf{1}$};

% === Projection arrows with tight midway labels ===
% Labels are anchored to the arrow midpoint with near-zero
% horizontal gap (inner sep = 2pt), eliminating the wide gap of
% the previous version.
\draw[arr=blue!60] (-3.0, 1.95) -- (-3.0, -0.05)
  node[midway, left=2pt, font=\footnotesize, blue!60,
       text width=2.1cm, align=right]
  {Gelfand\,+\,Connes\\ reconstr.\ (T-119)};

\draw[arr=green!50!black] (4.5, 1.95) -- (4.5, -0.05)
  node[midway, right=2pt, font=\footnotesize, green!50!black,
       text width=2.1cm, align=left]
  {Page--Wootters\\ clock algebra (T-118)};

% === M^4 level ===
\node[rectangle, rounded corners=8pt, draw=plospurple, fill=plospurple!10,
      minimum width=11cm, minimum height=1.3cm, align=center, font=\bfseries]
  (M4) at (0, -0.75)
  {$M^{4} = \mathbb{R}\times\Sigma^{3}$ \quad (T-120)};

\node[font=\footnotesize, color=plosgray] at (0, -1.85)
  {Product spectral triple $\to$ Einstein--Hilbert action
   $+$ Standard Model (T-121)};

% === Vacuum topology ===
\node[font=\footnotesize, color=plospurple] at (0, -2.45)
  {Vacuum: $\Sigma^{3}\cong S^{3}$ (de Sitter), $\Lambda > 0$
   \quad (T-120b)};

% === Key insight (bottom caption inside figure) ===
\node[font=\footnotesize\itshape, color=gray,
      text width=11cm, align=center]
  at (0, -3.30) {
  Spacetime is not fundamental. It is the macroscopic projection of
  the 7-dimensional coherence matrix onto the
  $\{A,S,D\}\!\cup\!\{O\}$ sector. Gravity is coherence seen from the
  external aspect.
};

\end{tikzpicture}

\caption{\textbf{Emergence of $M^4$ from the seven dimensions of $\Gam$.} The spatial sector
$\{A,S,D\}$ yields $\Sigma^3$ via Connes reconstruction (T-119). The temporal dimension O
yields $\mathbb{R}$ via Page--Wootters (T-118). The product spectral triple gives $M^4$
(T-120), from which the Einstein--Hilbert action follows (T-121).}
\label{fig:spacetime}
\end{figure}

\paragraph{Significance.}
Within UHM, spacetime is not the stage on which consciousness acts. Rather, the
7-dimensional coherence matrix $\Gam$ \emph{generates} spacetime as a macroscopic
projection (Figure~\ref{fig:spacetime}). The arrow of time
(the temporal modality $\triangleright$) is algebraically built into the subobject classifier
$\Omega$. Gravity is what coherence looks like from the external (physical) aspect.

\subsection{Consistency with quantum mechanics (sub-threshold limit)}
\label{subsec:qm-limit}

A critical requirement for any extension of physics: UHM must reduce to standard quantum
mechanics for systems that lack self-observation.

\begin{theorembox}[QM reduction for sub-threshold systems \status{T}]
For sub-threshold systems ($P < \Pcrit$), the viability gate $g_V(P) = 0$ shuts off
regeneration, and the dynamics reduce to the standard Lindblad master equation:
\begin{equation}
  \Lind_\Omega[\Gam]\big|_{P < \Pcrit} = -i[H, \Gam] + \mathcal{D}_0[\Gam]
  \label{eq:qm-limit}
\end{equation}
This is the dynamics of atoms, photons, and other elementary systems, which are unviable
($P < 2/7$) and cannot sustain self-reference.
\end{theorembox}

\paragraph{Significance.}
UHM is not a competitor to quantum mechanics but a proper extension: QM describes the
physics of systems without interiority; UHM describes autonomous systems where physical
and experiential aspects are inseparable. Every prediction of QM is preserved; UHM adds
predictions only for systems above the viability threshold.

\subsection{Resolution of the measurement problem}
\label{subsec:measurement}

The quantum measurement problem --- why does a superposition ``collapse'' to a definite
outcome, and in what basis? --- receives a three-part resolution in UHM:

\begin{enumerate}[leftmargin=2em]
  \item \textbf{Collapse as logical projection.} Wavefunction reduction is not a dynamical
        process but a \emph{logical operation}: projection of the state onto an atom of the
        subobject classifier $\Omega$. Each atom $S_k = |k\rangle\langle k|$ defines one
        outcome. The Born rule is encoded in the density matrix formalism:
        $p_k = \Tr(\chi_{S_k}\,\Gam)$~\status{T}. This is not an independent derivation
        but a consequence of the axioms (A1--A2), which assume the trace-probability link.
        What UHM adds is the identification of the measurement basis with the atoms of
        $\Omega$, resolving the preferred basis problem.
  \item \textbf{Preferred basis from Fano structure.} The ``preferred basis problem''
        (why does the system decohere in one basis rather than another?) is solved by
        the atomic structure of $\Omega$: the seven atoms
        $\{S_k\}_{k \in \{A,S,D,L,E,O,U\}}$ define the natural decomposition. By
        Zurek's einselection criterion (1981), the preferred basis consists of the
        fixed points of the decoherence channel --- which are precisely the
        Fano-structured atoms~\status{T}.
  \item \textbf{Self-measurement for conscious systems.} For systems with $R \geq 1/3$,
        measurement becomes \emph{self-observation}: the operator $\vp$ builds an internal
        model without irreversible information loss, unlike standard projective measurement
        which destroys coherences. Conscious measurement is a CPTP channel (information-preserving),
        not a projection (information-destroying)~\status{T}.
\end{enumerate}

Decoherence and measurement are thus dual aspects of the same mechanism: continuous
decoherence via $\mathcal{D}_\Omega$ in the limit $\gamma_k \to \infty$ contracts to
instantaneous projection onto an atom~\status{T}. The Lindblad operators
$L_k = \sqrt{\chi_{S_k}}$ are literally the square roots of logical truth values.
